\section{결론}\label{sec:conclusion}

SNN의 효율 약속은 이진 스파이크의 대수에 이미 들어 있다. 본 논문은
그 약속이 범용 MCU 배포 경로의 커널 경계에서 소멸한다는 것, 그리고
하드웨어 변경 없이 --- 세 가지 설계 원리(위치의 실체화, 불규칙
축의 인덱스화, 에필로그 불변)와 의미 보존 계약만으로 --- 비트
단위 동일하게 회수된다는 것을 STM32F746 실측으로 보였다. 결과는
종단 $6.12$배($2{,}328.5\rightarrow380.6$ms/이미지), 발화율 4.1\%
레이어의 동일 배치 dense 대비 13.0배이며, 여섯 빌드 변형의 출력
체크섬 일치가 정확성
주장을 실리콘에서 뒷받침한다. 표준 TinyML 스택은 그대로 남는다:
커스텀 연산자 등록만으로 통합되고 런타임 대가는 2.1\%다. 아울러 통념적
병목(명령 페치)에 대한 두 최적화가 실워크로드에서 무효과였다는
부정적 결과와 그 기제를 함께 보고했다 --- 병목의 위치는
마이크로벤치마크가 아니라 대상 워크로드 위의 절제 실측으로
확정되어야 한다.

향후 과제는 세 방향이다. 첫째, 최적화 후 병목이 된 고밀도
입력층이다: 인코더가 이벤트 목록을 직접 방출하게 하여 수용 영역
재스캔을 제거하는 것은 본 설계의 자연스러운 연장이고, 시간 코딩으로
입력 밀도 자체를 낮추는 것은 커널과 직교하는 지렛대다. 둘째, 비용
모델의 일반화다: 형상$\cdot$밀도가 다른 모델군과 다른 Cortex-M
등급에서 $\alpha,\beta$를 체계적으로 측정하면, 주어진 SNN에 대해
이벤트 커널의 채택 여부를 배포 전에 판정하는 도구가 된다. 셋째,
에너지 직접 계측과 정수 구현의 전 테스트셋 보드 정확도 실측으로
검증의 외연을 넓히는 것이다.
